\section{From an ultrapower embedding to a finite contradiction}
\label{sec:ultrapower-contradiction}
\label{sec:contradiction}

To prove Theorem~\ref{thm:main}, fix a nonzero unital nuclear
$C^*$-algebra $B$ and a free
ultrafilter $\omega$ on $\mathbb N$. A set of indices is called
$\omega$-large if it belongs to $\omega$.
Suppose, towards a contradiction,
that there is a unital injective $*$-homomorphism
\[
 \Phi:C^*(G)\longrightarrow B^\omega.
\]
We will associate a finite set of irreducible representations of
$\Gamma$ to each coordinate of $\Phi$.  Conjugation by $t$ will force
these sets to contain predecessors under restriction by $\theta$.
The resulting finite chains contradict the choice of $\pi_*$.

\subsection{Unitary lifts and a gap for varying matrix dimensions}

Write $u_g$ for the canonical unitary of $g\in G$ in $C^*(G)$.
Each unitary in $B^\omega$ has a representative consisting of unitaries
in $B$.  Indeed, if $(y_n)$ represents a unitary, then
\[
 \lim_\omega\|y_n^*y_n-1\|
 =\lim_\omega\|y_ny_n^*-1\|=0.
\]
On an $\omega$-large set both norms are less than $1/2$, so $y_n$ is
invertible and $y_n(y_n^*y_n)^{-1/2}$ is unitary.  The difference between
this unitary and $y_n$ tends to zero along $\omega$, by continuous
functional calculus.  Set the representative equal to $1$ outside this
set.

Choose such representatives $x_n(g)\in\mathcal U(B)$ of $\Phi(u_g)$
for every $g\in G$, with $x_n(e)=1$.  For each fixed pair $g,h\in G$,
\begin{equation}
 \lim_\omega\|x_n(g)x_n(h)-x_n(gh)\|=0.
 \label{eq:lift-defects}
\end{equation}
These choices impose no group relations at any individual coordinate.
All uses of \eqref{eq:lift-defects} below involve finitely many fixed
group elements; no intersection over all pairs is needed.
Put $\mathbf x_n=(x_n(s))_{s\in S}$.  Among the unitary equivalence
classes of nonzero finite-dimensional irreducible representations
of $\Gamma$, define
\begin{equation}
 D_n=\{[\pi]:d(\mathbf x_n,\pi)<a\}.
 \label{eq:coordinate-D}
\end{equation}
Here brackets denote unitary equivalence classes.  Each $D_n$ is finite
by Lemma~\ref{lem:nuclear-finiteness}, which applies to arbitrary unitary
$S$-tuples.

The next observation provides uniform control even when the dimensions
of the representations increase with $n$.

\begin{lemma}\label{lem:varying-gap}
For every sequence $(\pi_n)$ of nonzero finite-dimensional unitary
representations of $\Gamma$, with no bound on their dimensions,
\begin{equation}
 \lim_\omega d(\mathbf x_n,\pi_n)\in\{0\}\cup[\kappa,R].
 \label{eq:varying-gap}
\end{equation}
In particular, if
$\{n:d(\mathbf x_n,\pi_n)<a\}\in\omega$, then this limit is zero.
\end{lemma}

\begin{proof}
Let $C_n=B\otimes\mathcal B(\overline{V_{\pi_n}})$, and form the norm
ultraproduct
\[
 C_\omega=\prod_n C_n\big/
 \{(c_n):\lim_\omega\|c_n\|=0\}.
\]
The product consists of bounded sequences.  For every such sequence,
\begin{equation}
 \|[(c_n)]\|=\lim_\omega\|c_n\|.
 \label{eq:ultraproduct-norm}
\end{equation}
For the lower bound, adding an $\omega$-null sequence does not change
the ultralimit of the norms.  For the upper bound, if this ultralimit is
$c$, replace $c_n$ by $0$ outside
$\{n:\|c_n\|<c+\varepsilon\}\in\omega$ and then let
$\varepsilon\downarrow0$.

Equation~\eqref{eq:lift-defects} shows that
\[
 g\longmapsto[(x_n(g)\otimes\overline{\pi_n}(g))_n],
 \qquad g\in\Gamma,
\]
is a genuine unitary representation in the nonzero unital algebra
$C_\omega$.  By Lemma~\ref{lem:spectral-gap}, its Laplacian has
spectrum in $\{0\}\cup[\kappa,R]$.  Write
\[
 P_n=L\big((x_n(s)\otimes\overline{\pi_n}(s))_{s\in S}\big).
\]
Since $0\leq P_n\leq R1$, equation~\eqref{eq:ultraproduct-norm} gives
\begin{align*}
 \min\operatorname{spec}[(P_n)]
 &=R-\|R1-[(P_n)]\|\\
 &=\lim_\omega\bigl(R-\|R1-P_n\|\bigr)
 =\lim_\omega d(\mathbf x_n,\pi_n).
\end{align*}
This proves \eqref{eq:varying-gap}.  In the final assertion the limit
is at most $a<\kappa$, and hence must be zero.
\end{proof}

\subsection{The initial representation occurs}

We next use injectivity of $\Phi$ to ensure that $[\pi_*]\in D_n$
for an $\omega$-large set of indices.

\begin{lemma}\label{lem:initial-match}
We have $\lim_\omega d(\mathbf x_n,\pi_*)=0$, and consequently
$\{n:[\pi_*]\in D_n\}\in\omega$.
\end{lemma}

\begin{proof}
Induce $\pi_*$ from $\Gamma$ to $G$ on the Hilbert direct sum of
copies of $V_{\pi_*}$ indexed by the left cosets $G/\Gamma$.
Explicitly, choose representatives $q_i$, including the identity:
when $gq_i=q_j\gamma$ with $\gamma\in\Gamma$, the action sends
coordinate $i$ to coordinate $j$ by $\pi_*(\gamma)$.  This defines
a unitary representation.  The identity coordinate is invariant
under $\Gamma$ and carries $\pi_*$.  Thus the restriction of this representation
of $G$ contains $\pi_*$.  Tensoring with $\overline{\pi_*}$ gives a
nonzero $\Gamma$-invariant vector: on this coordinate it is
$\sum_j e_j\otimes\overline{e_j}$ for an orthonormal basis $(e_j)$ of
$V_{\pi_*}$.

Every unitary representation of $G$ gives a representation of the
full group algebra $C^*(G)$.  The preceding invariant vector shows that
\[
 \min\operatorname{spec}
 L\big((u_s\otimes\overline{\pi_*}(s))_{s\in S}\big)=0.
\]
Indeed, an invertible positive element could not acquire a kernel
under a unital representation.  The matrix amplification of $\Phi$
is injective and preserves spectra, so the corresponding spectral
minimum in $B^\omega\otimes\mathcal B(\overline{V_{\pi_*}})$ is also
zero.

For each fixed positive integer $d$ there is the canonical
identification
\[
 M_d(B^\omega)\cong (M_d(B))^\omega.
\]
To check its kernel, a sequence of $d$-by-$d$ matrices tends to zero in
norm along $\omega$ precisely when each of its finitely many entries
does.  Entrywise lifts give surjectivity.  Applying this identification
with $d=\dim\pi_*$ and the calculation of spectral minima in
Lemma~\ref{lem:varying-gap} proves the assertion.  This fixed matrix
amplification is distinct from the varying algebras used in that lemma.
\end{proof}

\subsection{A predecessor for every representation at one coordinate}

For representations $\alpha$ and $\beta$ of the same group, write
$\alpha\prec\beta$ when $\alpha$ is unitarily equivalent to a
subrepresentation of $\beta$.  The required quantifier order is the
following: on one $\omega$-large set, every member of $D_n$ has a
predecessor in that same finite set.

\begin{lemma}\label{lem:uniform-predecessor}
The set
\[
 \{n:(\forall[\pi]\in D_n)(\exists[\pi']\in D_n)
 \quad\pi\prec\pi'\circ\theta\}
\]
belongs to $\omega$.
\end{lemma}

\begin{proof}
Suppose otherwise.  On an $\omega$-large set $F$, choose an
irreducible $\pi_n$ with $[\pi_n]\in D_n$ for which no such predecessor
exists.  Outside $F$, choose the trivial representation.  By
Lemma~\ref{lem:varying-gap},
\begin{equation}
 \lim_\omega d(\mathbf x_n,\pi_n)=0.
 \label{eq:chosen-zero}
\end{equation}

Set $\alpha_n=\pi_n\circ\theta^{-1}$, a representation of $\Lambda$,
and let $\rho_n=\operatorname{Ind}_{\Lambda}^{\Gamma}\alpha_n$.
We will prove that $\lim_\omega d(\mathbf x_n,\rho_n)=0$ and that
every irreducible constituent $\pi'$ of $\rho_n$ satisfies
$\pi_n\prec\pi'\circ\theta$.
The block-minimum identity~\eqref{eq:block-minimum} will then select a
constituent belonging to $D_n$ on an $\omega$-large set, and the second
assertion will make that constituent a predecessor of $\pi_n$.

\emph{Transport to $\Lambda$.}
Represent $B$ faithfully and unitally on a Hilbert space $\mathcal H$,
and use the same notation for its unitaries on $\mathcal H$.  Choose
unit vectors $\xi_n\in\mathcal H\otimes\overline{V_{\pi_n}}$ such that
\[
 \sum_{s\in S}
 \|(x_n(s)\otimes\overline{\pi_n}(s))\xi_n-\xi_n\|^2
 <d(\mathbf x_n,\pi_n)+1/n.
\]
Because $\omega$ is free, $1/n$ tends to zero along $\omega$.  Thus
\begin{equation}
 \lim_\omega
 \|(x_n(g)\otimes\overline{\pi_n}(g))\xi_n-\xi_n\|=0
 \qquad(g\in\Gamma).
 \label{eq:fixed-word-invariance}
\end{equation}
Here is the fixed-word justification.  Put
$U_n(g)=x_n(g)\otimes\overline{\pi_n}(g)$.  For each $s\in S$,
$x_n(s^{-1})-x_n(s)^*$ tends to zero by
\eqref{eq:lift-defects}; a unitary and its inverse have the same
displacement norm on a vector.  For a fixed word
$g=s_1^{\varepsilon_1}\cdots s_k^{\varepsilon_k}$, repeated use of
\eqref{eq:lift-defects} replaces $U_n(g)$ by the corresponding product
of $U_n(s_j)^{\varepsilon_j}$ with an error tending to zero.
The displacement of that product is
at most the sum of the $k$ displacements of its factors.  Both the
word and $k$ are fixed while $n$ varies.

Put $\eta_n=(x_n(t)\otimes1)\xi_n$.  These are unit vectors.  For
$h\in\Lambda$ put $g=\theta^{-1}(h)$.  The identity $ht=tg$ and
\eqref{eq:lift-defects} give
\[
 \lim_\omega\|x_n(h)x_n(t)-x_n(t)x_n(g)\|=0.
\]
Together with \eqref{eq:fixed-word-invariance}, this proves
\begin{equation}
 \lim_\omega
 \|(x_n(h)\otimes\overline{\alpha_n}(h))\eta_n-\eta_n\|=0
 \qquad(h\in\Lambda).
 \label{eq:Lambda-invariance}
\end{equation}

\emph{Induction to $\Gamma$.}
We transfer these vectors from $\Lambda$ to $\Gamma$ by a fixed
eight-coordinate induction.  Use the left coset representatives
$r_1=e,r_2,\ldots,r_8$ of $\Gamma/\Lambda$ given by the translations
in $\{0,1\}^3$.  For $g\in\Gamma$ and $1\leq i\leq8$, define
$j(g,i)$ and $h(g,i)\in\Lambda$ uniquely by
\[
 gr_i=r_{j(g,i)}h(g,i).
\]
The induced representation $\rho_n$ acts on
$\bigoplus_{i=1}^8 V_{\pi_n}$: $\rho_n(g)$ sends coordinate
$i$ to coordinate $j(g,i)$ by $\alpha_n(h(g,i))$.  The relation
\[
 h(g_1g_2,i)=h(g_1,j(g_2,i))h(g_2,i)
\]
verifies the representation law.  Permutation of the coordinates and
unitarity of $\alpha_n$ verify unitarity.  Its conjugate has the same
coordinate rule with $\overline{\alpha_n}$ in place of $\alpha_n$.

Define $\zeta_n\in\mathcal H\otimes\overline{V_{\rho_n}}$ by
\[
 (\zeta_n)_i=8^{-1/2}(x_n(r_i)\otimes1)\eta_n.
\]
The eight summands are orthogonal and each has norm $8^{-1/2}$, so
$\|\zeta_n\|=1$.  Fix $s\in S$ and $i$, and write
$j=j(s,i)$ and $h=h(s,i)$.  The component sent from $i$ to $j$ by
$x_n(s)\otimes\overline{\rho_n}(s)$ differs from $(\zeta_n)_j$
by a vector whose norm is at most
\begin{align*}
 8^{-1/2}\big(&\|x_n(s)x_n(r_i)-x_n(r_j)x_n(h)\|\\
 &+\|(x_n(h)\otimes\overline{\alpha_n}(h))\eta_n-\eta_n\|\big).
\end{align*}
The first term tends to zero by \eqref{eq:lift-defects} and
$sr_i=r_jh$; the second does so by \eqref{eq:Lambda-invariance}.
There are only $8|S|$ fixed choices involved.  Hence
\[
 \lim_\omega\sum_{s\in S}
 \|(x_n(s)\otimes\overline{\rho_n}(s))\zeta_n-\zeta_n\|^2=0.
\]
Thus $\lim_\omega d(\mathbf x_n,\rho_n)=0$, and in particular
$d(\mathbf x_n,\rho_n)<a$ on an $\omega$-large set.
All error bounds are independent of $\dim\pi_n$.

\emph{Selecting a predecessor.}
Decompose the finite-dimensional unitary representation $\rho_n$ into
irreducible summands.  The corresponding Laplacian is a direct sum,
so its spectral minimum is the minimum of the spectral minima of
these summands.  Consequently, on an $\omega$-large set, some
irreducible constituent $\pi'_n$ of $\rho_n$ has $[\pi'_n]\in D_n$.

Every irreducible constituent $\pi'$ of $\rho_n$, including the one
just obtained, satisfies $\pi_n\prec\pi'\circ\theta$. This is the
finite-index form of Frobenius reciprocity. To verify it for every
constituent, let $W$ be its invariant subspace
and let $P_1:W\to V_{\pi_n}$ be projection onto the identity
coordinate.  This map is $\Lambda$-equivariant for the target
representation $\alpha_n$.  It is nonzero: if $w\in W$ has a nonzero
$i$th coordinate, then
\[
 P_1\rho_n(r_i^{-1})w=w_i\ne0.
\]
Since $\alpha_n$ is irreducible, $P_1|_W$ is surjective.  Its adjoint
is an injective $\Lambda$-intertwiner; after normalization by a
positive scalar it is an isometric intertwiner, because
$(P_1|_W)(P_1|_W)^*$ commutes with the irreducible $\alpha_n$.
It follows that $\alpha_n\prec\pi'|_\Lambda$; composing with
$\theta:\Gamma\to\Lambda$ gives the asserted containment.

Intersect the $\omega$-large set on which such a $[\pi'_n]$ belongs to
$D_n$ with $F$.  This contradicts the choice of $\pi_n$ and proves
the lemma.
\end{proof}

\subsection{Finite chains force periodicity}

The analytic argument is complete.  We finish with the elementary
consequence of finiteness that makes the predecessor property
incompatible with $\pi_*$. Consider a chain of finite-dimensional
irreducible unitary representations $\tau_0=\pi,\tau_1,\ldots$ of
$\Gamma$ whose classes lie in a
finite set $D$ and satisfy $\tau_i\prec\tau_{i+1}\circ\theta$.
Finiteness gives indices $l\geq0$ and $p\geq1$ with
$\tau_l\simeq\tau_{l+p}$.
Figure~\ref{fig:finite-chain} separates the initial segment of this
chain from the repeated segment.
The key point is that periodicity at the repeated class forces every
pullback of that representation to remain irreducible, so that
periodicity reaches the starting class.

\begin{figure}[htbp]
\centering
\begin{tikzpicture}[>=Stealth, every node/.style={font=\small}]
  \node (start) at (0,0) {$[\tau_0]=[\pi]$};
  \node (tail) at (2,0) {$\cdots$};
  \node (cycle) at (3.7,0) {$[\tau_l]$};
  \node (middle) at (5.4,0) {$\cdots$};
  \node (repeat) at (7.4,0) {$[\tau_{l+p}]$};
  \draw[->] (start) -- (tail);
  \draw[->] (tail) -- (cycle);
  \draw[->] (cycle) -- (middle);
  \draw[->] (middle) -- (repeat);
  \draw[dashed] (cycle.north) to[out=55,in=125]
     node[above] {same equivalence class} (repeat.north);
  \node at (1.7,-0.65) {initial segment};
  \node at (5.55,-0.65) {$p$ predecessor steps};
\end{tikzpicture}
\caption{A chain inside one finite set $D$. An arrow from $[\tau_i]$ to
$[\tau_{i+1}]$ means $\tau_i\prec\tau_{i+1}\circ\theta$.
Repetition gives $\tau_l\simeq\tau_l\circ\theta^p$ by equality of
dimensions. The proof then shows that each pullback of $\tau_l$ is
irreducible, allowing this periodicity to pass back to $\pi$.
The initial segment can be empty.}
\label{fig:finite-chain}
\end{figure}


\begin{lemma}\label{lem:finite-predecessor-cycle}
Let $\theta:\Gamma\to\Gamma$ be a homomorphism, and let $D$ be a
finite set of equivalence classes of nonzero finite-dimensional
irreducible unitary representations of $\Gamma$.  Suppose that every
$[\pi]\in D$ has $[\pi']\in D$ with $\pi\prec\pi'\circ\theta$.
Then for every $[\pi]\in D$ there is an integer $p\geq1$ such that
$\pi\simeq\pi\circ\theta^p$.
\end{lemma}

\begin{proof}
Starting from $\tau_0=\pi$, choose a chain in $D$ with
$\tau_i\prec\tau_{i+1}\circ\theta$.  Finiteness gives
$\tau_l\simeq\tau_{l+p}$ for some $l\geq0$ and $p\geq1$.
Composing the inclusions along this segment yields
\[
 \tau_l\prec\tau_{l+p}\circ\theta^p
 \simeq\tau_l\circ\theta^p.
\]
The first and last representations have the same dimension, so
\begin{equation}
 \tau_l\simeq\tau_l\circ\theta^p.
 \label{eq:cycle-periodicity}
\end{equation}
It follows by iteration that $\tau_l\circ\theta^{kp}\simeq\tau_l$
for every integer $k\geq0$.

For every $j\geq0$, the representation $\tau_l\circ\theta^j$ is
irreducible.  Indeed, if it had a nonzero proper invariant subspace,
that subspace would remain invariant for
$\tau_l\circ\theta^{kp}$ whenever $kp\geq j$, because
$\theta^{kp}(\Gamma)\subseteq\theta^j(\Gamma)$.  This contradicts
the irreducibility of $\tau_l\circ\theta^{kp}\simeq\tau_l$.
The initial segment of the chain gives
$\pi\prec\tau_l\circ\theta^l$.  Its right-hand side is irreducible,
so $\pi\simeq\tau_l\circ\theta^l$.  Composing
\eqref{eq:cycle-periodicity} with $\theta^l$ therefore gives
\[
 \pi\circ\theta^p
 \simeq\tau_l\circ\theta^{l+p}
 \simeq\tau_l\circ\theta^l
 \simeq\pi.
\]
\end{proof}

Choose one index $n$ in the intersection of the $\omega$-large sets
from Lemmas~\ref{lem:initial-match} and~\ref{lem:uniform-predecessor}.
Apply Lemma~\ref{lem:finite-predecessor-cycle} to $D_n$ and $\pi_*$.
It gives $\pi_*\simeq\pi_*\circ\theta^p$ for some $p\geq1$.
But $\pi_*$ acts nontrivially on the translations $\mathbb Z^3$,
whereas $\pi_*\circ\theta^p$ acts trivially on them: $\theta^p$
multiplies these translations by $2^p$, and $\pi_*$ factors through
reduction modulo $2$.  This is the required contradiction.

The argument used only that $B$ is nonzero, unital and nuclear, and
that $\omega$ is free.  The group $G$ and the algebra $C^*(G)$ are
independent of these choices.  Thus the same algebra has no unital
embedding into $B^\omega$ for any such $B$ and any free $\omega$.
Taking $B=\mathcal O_2$ proves the stated negative answer to
Kirchberg's embedding problem.
